% !TeX root = ../tesis.tex
%%%%%%%%%%%%%%%%%%%%%%%%%%%%%%%%%%%%%%%%%%%%%%%%%%%%%%%%%%%%%%%%%%%
% SECCIÓN 6.1 — EL CORREDOR ANILLAR COMO PROPUESTA
% Issue #25 (absorbe #14: definición formal del corredor anillar)
%%%%%%%%%%%%%%%%%%%%%%%%%%%%%%%%%%%%%%%%%%%%%%%%%%%%%%%%%%%%%%%%%%%
\section{El Corredor Anillar como Propuesta}
\label{sec:corredor-anillar}

El Capítulo~\ref{ch:analisis-red-actual} terminó con el Espectro MJG de la red
actual (§\ref{sec:diagnostico-garibelt}); ese perfil es la línea base de este
capítulo. Aquí se usa otra vez la cadena \texttt{Apímetro},
\texttt{VFTModel} y \texttt{Transport-gis-zmvm-mjg}, ahora sobre una red a la
que se le agrega un corredor anillar en el Anillo Periférico, y sus
indicadores se comparan con los de la red actual. Los casos son los del
glosario (§\ref{subsubsec:glosario-corredor-anillar}). El Caso~1 es la red
actual. El Caso~2, que en las figuras aparece como Anillo MB, suma el
corredor con la operación del Metrobús; el Caso~3, o Anillo Metro, lo suma con
la del Metro. Interesa saber si el método se puede repetir sobre una red
modificada y si sus indicadores captan la diferencia. La demanda, el costo y la
operación del corredor quedan fuera; lo que se reporta es cómo cambia la
estructura de la red con el anillo, bajo los supuestos de
§\ref{sec:anillo-escenarios}.

El capítulo avanza en cinco pasos. Esta sección define qué se entiende por
corredor anillar en la tesis, explica por qué se propone sobre el Anillo
Periférico y describe sus cuatro segmentos. La sección~\ref{sec:anillo-escenarios}
fija los parámetros de los Casos~2 y~3 y resume cómo se procesaron. Después,
la sección~\ref{sec:anillo-indicadores} revisa la cobertura, la Fuerza Capilar,
el tiempo de viaje y la intermediación de la red con anillo, siempre frente al
Caso~1. Con esos resultados, la sección~\ref{sec:anillo-espectro} arma el
Espectro MJG de los tres casos, y la sección~\ref{sec:anillo-interpretacion}
los lee desde el urbanismo.

\subsection{Definición Operativa del Corredor Anillar}
\label{subsec:corredor-anillar-definicion}

El glosario describe el corredor de alta capacidad anillar como una línea de
transporte público con infraestructura propia cuyo trazo rodea el área central
en lugar de converger hacia ella (§\ref{subsubsec:glosario-corredor-anillar}).
Para llevar esa idea al grafo hay que precisarla. En esta tesis se le llama
corredor anillar a una línea que cumple cuatro condiciones:

\begin{itemize}
    \item Su trazo forma un circuito cerrado alrededor del área central y no la
    atraviesa.
    \item Opera sobre vía exclusiva o confinada, de modo que el tráfico general
    la afecta poco o nada ($CF$ de 1.000 o 1.152,
    Cuadro~\ref{tab:cf-derecho-via}).
    \item Una parte de sus estaciones coincide con los cruces de las líneas
    radiales existentes, y ahí se generan los transbordos con el resto de la
    red.
    \item Su función es unir zonas periféricas entre sí, no llevar viajes al
    centro.
\end{itemize}

La diferencia con el corredor radial ya se planteó en el Capítulo~1 con los
esquemas de las Figuras~\ref{fig:transporte-radial}
y~\ref{fig:transporte-circular}. Aquí interesa en términos del grafo. En una
red radial, un viaje entre dos puntos de la periferia tiene que entrar hacia el
centro, cambiar de línea en alguno de los pocos nodos donde se cruzan las
radiales y volver a salir; por eso esos nodos acumulan intermediación. El
corredor anillar abre otra ruta posible: el viaje puede tomar el anillo en el
cruce más cercano y rodear el centro sin pasar por él. El
Cuadro~\ref{tab:radial-anillar} resume los rasgos que separan a los dos tipos de
corredor y el indicador de la tesis en el que cada rasgo debería notarse.
Rodear el centro tiene un costo que conviene reconocer desde ahora. Visto con
el Factor de Desviación ($DI$, §\ref{subsec:ruta-directa}), el anillo puede
alargar el recorrido sobre la red en viajes cuyo origen y destino están cerca
en línea recta, y en esos casos la red resulta menos eficiente en términos
geométricos. Lo que se busca con el anillo es otra cosa: descargar las rutas y
los nodos centrales por los que hoy pasa la mayor parte de los recorridos, y
dar alternativas de conexión entre las distintas zonas de la \zmvm{}.

\begin{table}[H]
    \centering
    \caption[Corredor radial y corredor anillar]{%
    Rasgos del corredor radial y del corredor anillar e indicador de la
    tesis que los registra.}
    \label{tab:radial-anillar}
    \small
    \begin{tabular}{|p{3.2cm}|p{4.2cm}|p{4.2cm}|p{2.4cm}|}
        \hline
        \textbf{Rasgo} & \textbf{Corredor radial} &
        \textbf{Corredor anillar} & \textbf{Indicador} \\
        \hline
        Geometría & Converge en el área central & Rodea el área central
            sin cruzarla & --- \\
        \hline
        Transbordos & Se concentran en pocos nodos centrales & Se reparten
            en los cruces con las radiales & $FC$, $B(v)$ \\
        \hline
        Viaje entre periferias & Pasa por el centro & Puede seguir el anillo &
            $T$ \\
        \hline
        Carga de los nodos centrales & Alta & Menor, si el anillo capta
            parte de los recorridos & $B(v)$ \\
        \hline
    \end{tabular}
    \fuentefigura{Elaboración propia.}
\end{table}

\subsection{Pertinencia en la \texorpdfstring{\zmvm{}}{ZMVM} y Localización sobre el Anillo Periférico}
\label{subsec:corredor-anillar-pertinencia}

La estructura que rodea la ciudad ya existe. El Anillo Periférico se construyó
como vía rápida para el automóvil (§\ref{subsec:sesgo-vial}) y hasta hoy no
tiene un carril confinado ni exclusivo para el transporte público
(§\ref{sec:justificacion}). Eso no significa que el transporte público no lo
use. Rutas de \textsc{rtp} y de Corredores Concesionados circulan por el
Periférico o lo recorren en parte de su trayecto; las rutas 57-A y 57-C de
\textsc{rtp}, entre Toreo y Metro Constitución de 1917, suman unos 83~km y
tienen un ramal sobre el Boulevard Manuel Ávila Camacho. Lo hacen, sin embargo,
en vía compartida con el tráfico general. El Coeficiente de Fricción Vial
recoge ese efecto (§\ref{subsec:friccion-vial}): con $CF = 1.380$, la
velocidad base de 11~km/h de ambos sistemas baja a unos 8.0~km/h
(§\ref{subsec:fase2-cf} y §\ref{subsec:validacion-impedancia}). Sobre el
Periférico, entonces, el transporte público avanza a la velocidad del
tráfico. Por esa razón la propuesta no consiste en sumar más rutas de
superficie, sino en un corredor de transporte masivo con vía propia, como el
Metro o el Metrobús, cuya velocidad dependa poco o nada de la congestión.

La pertinencia del anillo sale del diagnóstico de la red actual. En el
Capítulo~\ref{ch:analisis-red-actual} la intermediación aparece concentrada en
unos pocos nodos del eje Tacubaya--Mixcoac y del centro: Tacubaya
($B(v) = 0.2232$) y Mixcoac (0.1923) encabezan la jerarquía, y ocho de los
diez nodos con mayor $B(v)$ están en Benito Juárez o Cuauhtémoc
(§\ref{subsec:fase3-centralidad}). La Fuerza Capilar queda en banda Débil
(§\ref{subsec:fuerza-capilar}) y la cobertura cae al salir de la \cdmx{}: es de
53.52\,\% en la ciudad y de 3.92\,\% en los 60 municipios conurbados, 42 de los
cuales no tienen cobertura (§\ref{subsec:cobertura-espacial}). El Espectro MJG
resume ese cuadro con la Centralidad y la Accesibilidad en banda Crítica
(Cuadro~\ref{tab:espectro-mjg-diagnostico}). Es el patrón que el
Cuadro~\ref{tab:radial-anillar} atribuye a una red radial: los viajes entre
periferias pasan por el centro y los nodos centrales cargan con ellos. El
anillo se propone para abrir rutas que no dependan de esos nodos, sobre todo en
los contornos donde la red estructurada es más escasa.

El trazo propuesto sigue el Periférico en todo su recorrido y, por eso, no se
limita a la \cdmx{}. Como se observa en la Figura~\ref{fig:propuesta-anillo},
el anillo corresponde a lo que el glosario describe como un arco concéntrico
entre el área central y la primera corona periférica: rodea las alcaldías
centrales sin entrar en ellas y, en el norte y el oriente, sale de la ciudad.
De sus 94 estaciones, 66 quedan en trece alcaldías y 28 en cuatro municipios
del Estado de México: Naucalpan de Juárez, Tlalnepantla de Baz, Ecatepec de
Morelos y Nezahualcóyotl. Es, entonces, una propuesta de escala metropolitana,
que toca justo la franja donde la cobertura del transporte estructurado se
interrumpe. Los Casos~2 y~3 comparten ese trazo y esas estaciones; difieren
solo en la operación, que se detalla en §\ref{sec:anillo-escenarios}.

\begin{figure}[H]
    \centering
    \includegraphics[width=0.92\textwidth]{%
        Figures/Mapas/Qgis_Resources/RedConAnillo/Escenario_Metro/Base/Propuesta_Base_Anillar_Metro}
    \caption[Trazo del corredor anillar sobre el Anillo Periférico]{%
    Trazo del corredor anillar propuesto sobre el Anillo Periférico y red
    de transporte de la \zmvm{}. Los Casos~2 (Anillo MB) y~3 (Anillo Metro)
    comparten el trazo y las 94 estaciones. Corrida: 2026-10-04.}
    \label{fig:propuesta-anillo}
    \fuentefigura{Elaboración propia con \textbf{QGIS};
    datos: \texttt{Apímetro} y \texttt{VFTModel}~\citep{cabrera2025apimetro, cabrera2026VFTModel}.}
\end{figure}

\subsection{Segmentos y Cuadrantes del Anillo}
\label{subsec:corredor-anillar-cuadrantes}

El anillo se divide en cuatro segmentos, que en \texttt{Apímetro} se registran
como cuatro líneas: Periférico Interior Norte, Oriente, Sur y Poniente. Cada
una opera en los dos sentidos, de modo que la propuesta suma ocho ramales. Los
segmentos se unen en cuatro estaciones compartidas ---Parque Naucalli,
Colonias de Aragón, Periférico Oriente/Tláhuac y Luis Cabrera--- y así forman
un circuito cerrado. De las 94 estaciones, 20 se diseñaron como estaciones de
transbordo, contando las cuatro de unión; las demás son intermedias. El
Cuadro~\ref{tab:segmentos-anillo} resume cada segmento.

\begin{table}[H]
    \centering
    \caption[Segmentos del corredor anillar]{%
    Segmentos del corredor anillar propuesto: extremos, número de
    estaciones y demarcaciones que recorren.}
    \label{tab:segmentos-anillo}
    \small
    \begin{tabular}{|l|p{4.4cm}|c|p{5.4cm}|}
        \hline
        \textbf{Segmento} & \textbf{Extremos} & \textbf{Estaciones} &
        \textbf{Demarcaciones} \\
        \hline
        Norte & Parque Naucalli -- Colonias de Aragón & 24 &
            Naucalpan de Juárez, Tlalnepantla de Baz, Azcapotzalco,
            Gustavo A. Madero, Ecatepec de Morelos, Nezahualcóyotl \\
        \hline
        Oriente & Colonias de Aragón -- Periférico Oriente/Tláhuac & 18 &
            Nezahualcóyotl, Venustiano Carranza, Iztacalco, Iztapalapa \\
        \hline
        Sur & Periférico Oriente/Tláhuac -- Luis Cabrera & 25 &
            Iztapalapa, Tláhuac, Xochimilco, Tlalpan, Coyoacán,
            Álvaro Obregón, La Magdalena Contreras \\
        \hline
        Poniente & Luis Cabrera -- Parque Naucalli & 31 &
            La Magdalena Contreras, Álvaro Obregón, Benito Juárez,
            Miguel Hidalgo, Naucalpan de Juárez \\
        \hline
    \end{tabular}
    \fuentefigura{Elaboración propia con datos de
    \texttt{Apímetro}~\citep{cabrera2025apimetro}. Las estaciones de unión
    se cuentan en los dos segmentos que enlazan; sin repetirlas, el anillo
    tiene 94.}
\end{table}

Para leer el anillo por zonas conviene dividirlo en cuatro cuadrantes, como se
observa en la Figura~\ref{fig:propuesta-anillo-cuadrantes}. En el
\textbf{norponiente}, el trazo pasa por Naucalpan de Juárez y Tlalnepantla de
Baz, una franja que hoy atienden casi solo rutas de \textsc{rtp} y de
Corredores Concesionados; ahí el transbordo previsto con la red estructurada
es Cuatro Caminos, terminal de la Línea~2 del Metro. En el \textbf{nororiente}
cruza Gustavo A. Madero, Ecatepec de Morelos y Nezahualcóyotl, y en Río de los
Remedios se encuentra con la Línea~B del Metro y con el Metrobús. El
\textbf{suroriente}, en Iztapalapa y Tláhuac, concentra los cruces con la
Línea~12 en Periférico Oriente, con la Línea~8 y el Cablebús en Constitución
de 1917 y con la Línea~A en Canal de San Juan. En el \textbf{surponiente}, el
anillo toca el Metrobús en Periférico Sur y la Línea~7 en Barranca del Muerto
y San Pedro de los Pinos, y une Coyoacán, Tlalpan, Álvaro Obregón y La
Magdalena Contreras sin pasar por el centro. Lo que cada cuadrante gana o
pierde con el anillo se revisa en §\ref{sec:anillo-indicadores}.
% La Figura fig:propuesta-anillo-cuadrantes se trasladó a §6.3.1
% (6-3-Indicadores-Escenario/6-3-1-Escenarios.tex, Caso 3) el 2026-10-18. Al redactar §6.3.1,
% revisar la frase «como se observa en la Figura…» de este párrafo
% (ahora es una referencia hacia adelante).

\subsection{Limitaciones del Modelo}
\label{subsec:corredor-anillar-limitaciones}

Con la definición, la justificación y el trazo, la propuesta ya se puede
llevar al grafo. Antes de hacerlo conviene dejar claro qué representa el modelo
y qué deja fuera. Las principales limitaciones de la propuesta son cuatro:

\begin{enumerate}
    \item \textbf{No incluye las condiciones físicas del trazo.} Los escenarios
    no consideran la pendiente del terreno, el ancho de las vialidades ni el
    tipo de calle o avenida que cruza cada segmento, y tampoco si el corredor
    iría elevado, bajo tierra o en un carril confinado del Periférico. Al grafo
    le basta saber dónde están las estaciones, con qué se conectan y a qué
    velocidad opera cada caso. Se trata, por lo tanto, de un modelo
    teórico-urbano-computacional que sirve de apoyo a la toma de decisiones,
    no de un proyecto ejecutivo.

    \item \textbf{No incluye la dimensión social.} La propuesta no incorpora
    datos de población, ingreso o demanda de viajes; la extensión del estudio
    no permitió llegar a ese análisis. Es una dimensión necesaria para evaluar
    cualquier corredor, y su estudio queda como trabajo futuro
    (Capítulo~\ref{ch:conclusiones}).

    \item \textbf{La división en cuatro segmentos es una decisión del autor.}
    En los dos casos el anillo se divide en Norte, Oriente, Sur y Poniente
    porque cada zona tiene condiciones geográficas, sociales y económicas
    distintas y plantea preguntas de análisis diferentes. Por la misma razón,
    la propuesta se denomina \textbf{Anillo Periférico Interior}
    (\textsc{api-i}) y sus segmentos, Periférico Interior Norte, Oriente, Sur y
    Poniente. Cuando se trazó el Periférico, durante la regencia de Uruchurtu
    (§\ref{sec:anillo-periferico}), rodeaba la ciudad de esa época; desde
    entonces la \zmvm{} creció de forma dispersa, con el patrón saltón
    descrito en la Figura~\ref{fig:sprawl}, y hoy el anillo ya no coincide
    con la periferia metropolitana. Lo que sí delimita es el borde exterior
    del Primer y el Segundo Contorno (Figura~\ref{fig:semovi-contornos}).

    \item \textbf{No aborda un anillo exterior.} De lo anterior se desprende la
    necesidad de estudiar un Anillo Periférico Exterior, que atienda la
    periferia actual de la \zmvm{}. Ese caso no se desarrolla en este
    documento: por la metodología que exige, la extensión del territorio y el
    tiempo disponible, requiere un estudio propio.
\end{enumerate}

Con estas limitaciones a la vista, §\ref{sec:anillo-escenarios} describe cómo
se incorporan los dos casos al grafo y con qué parámetros.
